\documentclass[12pt]{article}
\usepackage{amsmath, amssymb}
\usepackage[margin=1in]{geometry}
\setlength{\parskip}{6pt}
\title{QUBO Formulation: Clinical Enzyme Target Identification}
\date{}
\begin{document}
\maketitle

\subsection*{Clinical Enzyme Target Identification}

\paragraph{Problem Statement.}
Given a metabolic network containing $N$ candidate enzyme targets, each enzyme $i$ is associated with a benefit score $B_i$ quantifying its therapeutic impact, and a pairwise redundancy penalty $R_{ij}$ quantifying the pathway overlap between enzymes $i$ and $j$. The objective is to select exactly $K$ enzymes such that the total benefit is maximized while accounting for redundancy penalties among the selected pairs.

\paragraph{Decision Variables.}
For each enzyme candidate $i \in \{1, \dots, N\}$, we introduce a binary decision variable $x_i \in \{0, 1\}$. The variable $x_i = 1$ indicates that enzyme $i$ is selected as a target, and $x_i = 0$ otherwise.

\paragraph{QUBO Derivation.}
\textbf{Step 1.} The original maximization objective is given by
\begin{equation}
  \max_{\mathbf{x} \in \{0,1\}^N} \sum_{i=1}^N B_i x_i - \sum_{1 \le i < j \le N} R_{ij} x_i x_j.
\end{equation}
Since standard QUBO solvers minimize energy, we negate the objective to obtain the minimization form:
\begin{equation}
  \min_{\mathbf{x} \in \{0,1\}^N} \left( -\sum_{i=1}^N B_i x_i + \sum_{1 \le i < j \le N} R_{ij} x_i x_j \right).
\end{equation}

\textbf{Step 2.} The cardinality constraint requires exactly $K$ enzymes to be selected:
\begin{equation}
  \sum_{i=1}^N x_i = K.
\end{equation}

\textbf{Step 3.} We enforce the constraint via a quadratic penalty term with weight $\lambda > 0$:
\begin{equation}
  P(\mathbf{x}) = \lambda \left( \sum_{i=1}^N x_i - K \right)^2.
\end{equation}
Expanding the square and applying the idempotent property $x_i^2 = x_i$ for binary variables yields:
\begin{align}
  P(\mathbf{x}) &= \lambda \left( \sum_{i=1}^N x_i^2 - 2K \sum_{i=1}^N x_i + K^2 \right) \\
  &= \lambda \left( \sum_{i=1}^N x_i - 2K \sum_{i=1}^N x_i + K^2 \right) \\
  &= \lambda \left( (1 - 2K) \sum_{i=1}^N x_i + K^2 \right) \\
  &= \sum_{i=1}^N \lambda(1 - 2K) x_i + \lambda K^2.
\end{align}
The resulting general penalty contribution to the Hamiltonian is $\sum_{i=1}^N \lambda(1 - 2K) x_i + \lambda K^2$.

\textbf{Step 4.} Combining the negated objective and the penalty term gives the total QUBO Hamiltonian $H(\mathbf{x})$:
\begin{equation}
  H(\mathbf{x}) = \sum_{i=1}^N \left( -B_i + \lambda(1 - 2K) \right) x_i + \sum_{1 \le i < j \le N} R_{ij} x_i x_j + \lambda K^2.
\end{equation}

\textbf{Step 5.} Reading off the coefficients from the standard QUBO form $H(\mathbf{x}) = \mathbf{x}^\top Q \mathbf{x} + c$, we obtain the general closed-form expressions for the upper-triangular $Q$ matrix entries and the constant offset:
\begin{align}
  Q_{i,i} &= -B_i + \lambda(1 - 2K), \quad \forall i \in \{1, \dots, N\}, \\
  Q_{i,j} &= R_{ij}, \quad \forall 1 \le i < j \le N, \\
  c &= \lambda K^2.
\end{align}
By symmetry, $Q_{j,i} = Q_{i,j}$ for $i \neq j$.

\paragraph{Penalty Weight Justification.}
Let $W = \max_{i} |B_i|$ denote the maximum absolute benefit score. The maximum possible change in the objective function resulting from flipping a single variable is bounded by $W$. Therefore, setting $\lambda > W$ ensures that the penalty incurred by violating the cardinality constraint strictly dominates any potential gain from the original objective terms. This guarantees that any infeasible configuration has strictly higher energy than the optimal feasible configuration.

\paragraph{Correctness Argument.}
Minimizing $H(\mathbf{x})$ over $\{0,1\}^N$ recovers the optimal solution to the original problem because (i) feasible solutions satisfying $\sum_{i=1}^N x_i = K$ incur zero penalty and their energy exactly matches the negated objective value, and (ii) infeasible solutions violate the cardinality constraint, incurring a penalty of at least $\lambda$. Since $\lambda$ is chosen to strictly dominate the maximum possible variation in the objective terms, every infeasible solution has strictly higher energy than the optimal feasible solution. Consequently, the global minimum of $H(\mathbf{x})$ must correspond to a feasible assignment that minimizes the negated objective, which is equivalent to maximizing the original objective.

\paragraph{Illustrative Example.}
Consider a small instance with $N = 3$ enzymes, target cardinality $K = 2$, benefit scores $B_1 = 5$, $B_2 = 4$, $B_3 = 3$, and redundancy penalties $R_{12} = 1$, $R_{13} = 2$, $R_{23} = 1$. We set the penalty weight to $\lambda = 10$, satisfying $\lambda > \max_i |B_i|$. Instantiating the general formulas yields the concrete components:
\begin{align}
  \text{Negated objective:} & \quad -5x_1 - 4x_2 - 3x_3 + x_1x_2 + 2x_1x_3 + x_2x_3, \\
  \text{Penalty term:} & \quad 10(x_1 + x_2 + x_3 - 2)^2 = -30x_1 - 30x_2 - 30x_3 + 40.
\end{align}
Combining these gives the resulting Hamiltonian:
\begin{equation}
  H(\mathbf{x}) = -35x_1 - 34x_2 - 33x_3 + x_1x_2 + 2x_1x_3 + x_2x_3 + 40.
\end{equation}
Enumerating all $\binom{3}{2} = 3$ feasible assignments ($x_1+x_2+x_3=2$):
\begin{itemize}
  \item $\mathbf{x} = (1,1,0)$: $H = -35 - 34 + 1 + 40 = -28$.
  \item $\mathbf{x} = (1,0,1)$: $H = -35 - 33 + 2 + 40 = -26$.
  \item $\mathbf{x} = (0,1,1)$: $H = -34 - 33 + 1 + 40 = -26$.
\end{itemize}
The minimum energy is $-28$, achieved at $\mathbf{x}^* = (1,1,0)$. This corresponds to selecting enzymes 1 and 2, yielding a maximum total score of $5 + 4 - 1 = 8$.

\end{document}